%=============================================================================
% MeshCrowd: Trustworthy Crowd Sensing through Physical-Layer Security
% and Gamified Integrity Engineering
% IEEE two-column manuscript — cybersecurity & gamification focus.
% Compile with:  pdflatex MeshCrowd_Cyber.tex  (run twice for refs)
% Requires: IEEEtran class and the 'cite' package.
%=============================================================================
\documentclass[journal,10pt]{IEEEtran}

\usepackage{cite}
\usepackage{graphicx}
\usepackage{amsmath,amssymb}
\usepackage{textcomp}
\usepackage{array}
\usepackage{booktabs}
\usepackage{url}
\usepackage[hidelinks]{hyperref}
\usepackage{algorithm}
\usepackage{algpseudocode}

\graphicspath{{figures/}}

\begin{document}

%---------------------------------------------------------------------------
% Title block
%---------------------------------------------------------------------------
\title{MeshCrowd: Trustworthy Crowd Sensing through Physical-Layer Security and Gamified Integrity Engineering}

\author{%
\IEEEauthorblockN{Prepared for submission to IEEE --- draft}
\IEEEauthorblockA{Device-Free Passive Sensing on Commodity Phones: Secure Multi-Band Fusion, dBm-Based Ranging, and a Reliability Bonus that Rewards Honesty}
}

\markboth{Journal of \LaTeX\ Class Files,~Vol.~XX, No.~X, 2026}%
{Author \MakeLowercase{\textit{et al.}}: Trustworthy Crowd Sensing}

\maketitle

%---------------------------------------------------------------------------
% Abstract
%---------------------------------------------------------------------------
\begin{abstract}
Crowd-sourced device-free sensing is attractive precisely because it needs no
installed infrastructure: people already carry the radios, and the observed
target carries nothing. Yet the same openness that makes it cheap makes it
exploitable. A single malicious participant can inject forged beacons
(spoofed SSIDs, rogue access points, phantom Bluetooth advertising) or publish
forged crossing reports, and in a system that rewards reliability the mesh mean
--- the very quantity being rewarded --- becomes the attack surface. This paper
recasts \emph{MeshCrowd}, a device-free, multi-band passive RF sensing
system that runs entirely on an unmodified commodity smartphone, as a study in
two problems that are usually treated apart. The first is \emph{cyber}: how to
make a sensing mesh resilient to Byzantine contributors that report data they
never observed. The second is \emph{gamification}: how to design the incentives
so that the rational strategy for a participant is precisely the behaviour the
system wants, rather than merely a behaviour it can measure. We show that the
two are not independent. By computing an explicit estimated accuracy and
granting a reliable-accuracy bonus to the devices whose measurements lie
nearest the fused mesh mean, the incentive and the security objective coincide:
encouraging accuracy is equivalent to encouraging honesty. The system fuses
per-link dip depth, duration, temporal texture, variance, prominence and
multi-link coherence across BLE and dual-band (2.4~GHz and 5~GHz) Wi-Fi into
pedestrian/vehicle typing; converts measured dBm into a gated radial range band
$f \in [0.15, 0.92]$ with a hard refusal below the $-90$~dBm noise floor;
detects field change across $\approx$10~s scan epochs from the median of the
last 20 samples; and realises a two-node roadside transit counter over Classic
Bluetooth RFCOMM and BLE L2CAP. A mesh of 2--5 instances extends the useful
radius from 4--12~m toward 20--40~m with zero false confirmations across 160
empty scenes. We frame a physical-layer fingerprinting defence and a
reputation-weighted fusion rule as future work that converts the reliability
bonus from a social incentive into a cryptographic-grade trust primitive.
\end{abstract}

\begin{IEEEkeywords}
device-free sensing, physical-layer security, radio-frequency fingerprinting,
Byzantine resilience, mobile crowdsensing, gamification, reputation, reward
design, Bluetooth Low Energy, Wi-Fi sensing, received signal strength, change
detection, integrity.
\end{IEEEkeywords}

%---------------------------------------------------------------------------
% I. INTRODUCTION
%---------------------------------------------------------------------------
\section{Introduction}

\subsection{The Opportunity and the Adversary}
Passive (device-free) RF sensing observes how physical objects modify the
propagation of wireless signals. When a person or vehicle crosses a radio
link, it attenuates, scatters and diffracts the field, and the received signal
strength (RSSI) dips and recovers with a characteristic signature. Because the
target carries no device and the observer needs no installed hardware, this
paradigm underpins indoor localization, ambient-assisted living and occupancy
analytics~\cite{wilson2010,januszkiewicz2018}. A modern smartphone already
contains the required radios: a Bluetooth Low Energy (BLE) transceiver and a
dual-band (2.4~GHz and 5~GHz) Wi-Fi interface can observe the signal strength
of every advertising beacon and access point (AP) within range \emph{without}
connecting to it. Each observed link therefore becomes a measurement ray whose
attenuation state is monitored in real time~\cite{ramirez2021,milano2024}.

The economic and social argument is even stronger. Dedicated sensing
infrastructure is costly, and camera-based alternatives raise privacy concerns
in homes, workplaces and public spaces~\cite{khan2023,zhou2020,dmip2023}.
Smartphone-based passive sensing can be deployed for free by people who already
carry the device, and when many phones in the same area cooperate --- a
\emph{MeshCrowd} mesh --- the geometric diversity of their rays increases
both accuracy and spatial reach~\cite{li2024}.

This openness, however, is also the vulnerability. A device-free mesh has no
notion of who is contributing or whether a contribution is genuine. The observer
cannot authenticate a beacon it has not joined, and the target cannot be
fingerprinted because it carries nothing. A single adversarial participant can
forge beacons --- spoofed SSIDs, rogue APs, phantom BLE advertising --- or
publish crossing reports it never observed. In a system that computes its own
reputation from contributed data, the mesh mean becomes a target, and every
reward derived from it inherits the deception~\cite{zhang2022,wang2018}.

\subsection{Problem Statement}
Three difficulties compound. First, detection is physics-limited: occlusion
attenuation decays with distance through Fresnel-zone occupancy while the
environment's own fluctuation grows with range, so beyond some distance the
event is statistically indistinguishable from noise, and any system that does
not quantify this limit misleads its user. Second, RSSI ranging is notoriously
ambiguous: a log-distance model converts decibels into metres, but on weak
links the inverse function is extremely ill-conditioned and an unconstrained
estimate can be wrong by tens of metres~\cite{ramirez2021,milano2024}. Third,
and most relevant here, community mode introduces an \emph{integrity} problem:
forged beacons and forged reports both attack the mesh mean, and the
gamification layer that sustains participation doubles as the incentive to
deceive. The core question is whether the reward mechanism can be designed such
that reporting truthfully is the dominant strategy.

\subsection{Contribution}
This paper repositions \emph{MeshCrowd} as a case study joining
\emph{physical-layer security} and \emph{gamified incentive design}. The
contributions are as follows.
\begin{itemize}
  \item A \textbf{multi-band, device-free fusion engine} running on a commodity
  phone, fusing per-link dip depth, duration, temporal texture, variance,
  prominence and multi-link coherence across BLE and dual-band Wi-Fi into
  pedestrian/vehicle typing (1,512 Monte-Carlo scenes; BLE 57.7\,\% vs.\
  25.7\,\% and 16.1\,\% for the companion bands).
  \item A \textbf{gamified integrity engine}: an explicit estimated accuracy and
  a reliable-accuracy bonus granted to devices nearest the fused mesh mean,
  deliberately aligning the social incentive with the security objective.
  \item A \textbf{trustworthiness analysis} of the mesh under Byzantine
  contribution, framing a physical-layer fingerprinting defence and a
  reputation-weighted fusion rule as future work.
  \item An \textbf{honest ranging} component with an explicit refusal semantics:
  log-distance inversion, a quality ladder from $-20$~dBm to the $-90$~dBm
  noise floor, a gate at $-75$~dBm and a radial band $f \in [0.15, 0.92]$,
  validated on 288 scenes where the gate suppresses 14--23~m illusions.
\end{itemize}

%---------------------------------------------------------------------------
% II. RELATED WORK (cyber + gamification anchored)
%---------------------------------------------------------------------------
\section{Related Work}

\subsection{Physical-layer security and device fingerprinting}
Conventional wireless authentication is cryptographic, but the overhead is often
prohibitive for the low-cost, resource-constrained devices that populate the
Internet of Things (IoT). Physical-layer-based device fingerprinting offers a
passive alternative that is readily applicable to legacy hardware, and a
comprehensive survey distinguishes hardware-impairment-based radio-frequency
fingerprint identification (RFFI) from channel-feature-based (location)
authentication~\cite{zhang2025survey}. Multi-domain physical-layer fingerprint
features --- time- and frequency-domain plus modulation-domain statistics of the
frame preamble --- have been used to identify Wi-Fi IoT devices, reaching 98\,\%
for 15 device types and 90.76\,\% for 10 low-differentiation network
cards~\cite{zhang2023wifi}. Nonlinear signal analysis, such as detrended
fluctuation analysis (DFA) fused with carrier-frequency-offset and fractal-dimension
features, further strengthens authentication in dynamic channels without added
cryptographic overhead~\cite{dfa2025}. Critically, these fingerprints arise from
intrinsic manufacturing imperfections and are therefore hard to replicate,
which supplies the anti-spoofing property our defence requires. On the attack
side, undetectable attacks against physical-layer fingerprint authentication are
known, and integrated semi-supervised and environment-semantics-aware schemes
have been proposed to harden them~\cite{phyfinatt,esaNet}. We draw on this
literature to argue that a beacon trusted on the basis of identifiers alone is
insufficient, and that the mesh must instead anchor trust in physical-layer
properties and reputation.

\subsection{Crowd-sensing incentives and reputation}
Mobile-crowdsensing research is unambiguous that participation must be
engineered, and that quality-weighted rewards outperform raw volume.
Reputation-aware and privacy-preserving incentive schemes based on deep
reinforcement learning and Stackelberg games couple reward to trustworthiness
and preserve user privacy~\cite{zhang2022,wang2018}. These works establish that
reward should be a function of the \emph{quality} of the contribution, not its
\emph{quantity}, which is precisely the design choice our reliability bonus
makes. We extend this line by arguing that when the fusion estimate is itself
computed from the rewarded contributions, the incentive mechanism becomes a
defence primitive.

\subsection{Gamification for behaviour and security outcomes}
Gamification is established as a mechanism for behaviour change, and its
application to security is well studied. A systematic review found that
gamified cybersecurity initiatives outperform non-gamified formats on
engagement, motivation, awareness, skill acquisition, behaviour and retention,
while also noting the inconsistent use of controls~\cite{gwenhure2025}.
Gamified security e-training has been shown to improve employees' security
self-efficacy and, critically, to reduce actual phishing-click rates in a
controlled campaign~\cite{horst2024gamify}. Gamification applied to privacy and
information-security awareness in the public sector was perceived as more
motivating and effective for retention, with organisational culture and
leadership support as the main moderators~\cite{cm2026gamify}. Collectively,
these results validate the premise that a points-and-badges layer can drive
participation; our contribution is to make the reward \emph{aligned} with the
measurement so that honest reporting is the rational strategy.

\subsection{Counting, ranging and the sensing physics}
Device-free people counting exploits CSI and Doppler to estimate flow
\cite{khan2023,zhou2020,dmip2023}, and RSSI ranging requires integrity
monitoring against unstable beacons and obstacle-aware treatment of missing
links~\cite{ramirez2021,milano2024,bai2020}. Human-body shadowing at 2.4~GHz is
characterised and compensable~\cite{januszkiewicz2018}. Fisher-information
analysis bounds device-free localization in integrated sensing-and-communication
networks~\cite{li2024}, which supports our decision to report bands rather than
points at long range.

%---------------------------------------------------------------------------
% III. SYSTEM DESIGN
%---------------------------------------------------------------------------
\section{System Design}

\subsection{Sensing model and bands}
MeshCrowd observes BLE advertising beacons (2.4~GHz) and Wi-Fi APs at 2.4
and 5~GHz. Received power is modelled by log-distance path loss,
$\mathrm{RSSI}(d) = \mathrm{PL}_0 - 10\,n\,\log_{10}(d)$, with per-band
reference parameters $(n, R_{1\,\mathrm{m}}) = (2.4, -42)$~dBm, $(2.6, -38)$~dBm
and $(2.9, -35)$~dBm. Fluctuation is lognormal scintillation growing with range
(0.45--0.70~dB at short range, $\approx$0.055~dB/m). An occlusion is a
Gaussian-profiled attenuation $A_{\mathrm{eff}}\exp(-\mathrm{lat}^2/2\sigma^2)$
with amplitude scaled by a Fresnel term $(1+d/10)^{-0.6}$ and lateral width
$\sigma = 0.35$~m (pedestrian) or $1.15$~m (vehicle); vehicles scale the
occlusion to $\approx$25~dB with a 5.5~Hz flutter, and pedestrians carry a
3.1~Hz body wobble.

\subsection{Signal-processing chain}
Each link is a stream of (timestamp, dBm) samples processed by production
Kotlin classes: \texttt{RingWindow}, \texttt{Stats} (median, robust
dispersion), \texttt{BaselineTracker} (adaptive no-event level) and
\texttt{DipDetector}/\texttt{PeakDetector}. Three decisions emerged from
measured failures: a \emph{noise-aware trigger floor}, a \emph{sustained-trigger
gate}, and a \emph{stable-level policy}. Fig.~\ref{fig:method} shows the
functional decomposition and the artefacts each stage emits.

\begin{figure*}[t]
  \centering
  \includegraphics[width=\textwidth]{fig1_method.png}
  \caption{Functional decomposition of the sensing pipeline. Each stage
  produces artefacts --- per-link dBm tables, epoch tables, transit counts and
  mesh ledgers --- that the evaluation and the integrity analysis replay.}
  \label{fig:method}
\end{figure*}

\subsection{Features and fusion}
For each event the engine computes dip depth, duration, texture (jitter ratio),
variance, prominence, and coherence --- the number of links registering the
event in a common interval, which becomes the mesh primitive.
\texttt{EntityClassifier} (v1.1) applies a learned weight vector per band and
fuses at the crossing level. BLE carries typing, Wi-Fi 5~GHz contributes the
strongest attenuation ($A_0 = 11.5$~dB) and Wi-Fi 2.4~GHz adds coverage; the
bands are fused asymmetrically by design, and equal-weight fusion was rejected
because it dragged BLE's 57.7\,\% to 31\,\%.

\subsection{Ranging and the refusal semantics}
Ranging inverts the log-distance model with an explicit honesty contract. A
quality ladder runs from strong links ($-20$~dBm) to the $-90$~dBm noise floor;
estimation is gated at $-75$~dBm, below which an unconstrained estimate can be
off by 14--23~m. The system therefore publishes a radial band
$f \in [0.15, 0.92]$ rather than a point, and excludes links at or below
$-90$~dBm from range and accuracy mathematics, retaining them only for
multi-link agreement. The refusal is itself a security property: it prevents a
noisy or adversarial link from manufacturing a confident false range.

\begin{figure*}[t]
  \centering
  \includegraphics[width=\textwidth]{fig2_arch.png}
  \caption{Layered architecture. The inference and application layers, and
  the gamification layer, are the surfaces where trustworthiness must be
  enforced.}
  \label{fig:arch}
\end{figure*}

\subsection{Scan-epoch change detection}
The engine treats the field as two timed snapshots on a $\approx$10~s
moving-object time scale. Each epoch stores a timed record (timestamp, place,
geo, per-link dBm table); the engine resets and rescans; and the difference is
computed from stable levels --- the median of the last 20 samples --- or a
transient dip of $\geq 3$~dB. Single-sample comparison has $\approx$1~dB
standard deviation and produced a 25--62\,\% false-positive rate, whereas the
20-sample median collapses the fluctuation before any comparison is made.

\subsection{Two-node transit counter}
Two phones form a barrier across a passage or road using Classic Bluetooth
RFCOMM and BLE L2CAP. A hybrid link-presence/loss classifier treats both a link
drop and a link appearance as evidence and types the entity between the two
nodes, realising barrier-coverage geometry without synchronised nodes.

\subsection{Mesh agreement and the gamified reliability bonus}
Nearby instances exchange crossing records. A crossing is published only when
confirmed by a threshold number of nodes, and rays that disagree with the
majority are down-weighted, so the mesh mean becomes the fused estimate. The
gamification layer computes an \emph{explicit estimated accuracy} and grants a
\emph{reliability bonus} to the devices whose measurements lie nearest that
mean, plus a fusion multiplier when a crossing is confirmed by several nodes.
The alignment is deliberate and is the crux of this paper: because the reward
is computed from the same fused estimate that is published, encouraging
accuracy is identical to encouraging honest measurement.

\subsection{Deployment and gamification surface}
Fig.~\ref{fig:sense}--Fig.~\ref{fig:club} render the deployed build. The
\emph{Sense} tab is the live instrument. The \emph{Club} tab is where the
integrity mechanism is exposed to the participant: it reports mesh contribution
(nodes in your mesh, fused accuracy, crowd level) and the scoring rule, in which
people and vehicles earn XP and a mesh-confirmed crossing earns a $\times 2$
fusion bonus --- the social face of the alignment argument of
Section~\ref{sec:method}. The \emph{Settings} tab exposes every shipped constant
(including the range gate and the noise floor) because each was the outcome of a
calibration-and-rejection protocol, and the \emph{Privacy} panel commits to no
camera, no microphone and no stored beacon identifiers, which is the precondition
for a crowd-sensing system to be accepted at all.

\begin{figure}[t]
  \centering
  \includegraphics[width=0.58\columnwidth]{fig3_sense.png}
  \caption{The \emph{Sense} dashboard: 33 passive links, live counters
  (People 67, Vehicles 66, Disturbances 165, RF Links 33), moving variance
  11.45~dB$^2$, jitter 1.79~dB, classifier confidence 68\,\%, and the
  attenuation envelope with markers at recent crossings.}
  \label{fig:sense}
\end{figure}

\begin{figure}[t]
  \centering
  \includegraphics[width=0.58\columnwidth]{fig4_scans.png}
  \caption{The \emph{Scans} tab: a stored epoch and the verdict ``field
  CHANGED since previous scan'' over a 120-scan history.}
  \label{fig:scans}
\end{figure}

\begin{figure}[t]
  \centering
  \includegraphics[width=0.58\columnwidth]{fig5_mesh.png}
  \caption{The \emph{Mesh} tab: 2--5 cooperating nodes confirm each other's
  crossings, extending the radius from 4--12~m toward 20--40~m, with a
  reliability bonus for devices nearest the fused estimate.}
  \label{fig:mesh}
\end{figure}

\begin{figure}[t]
  \centering
  \includegraphics[width=0.58\columnwidth]{fig6_transit.png}
  \caption{The \emph{Transit} tab: a two-node barrier over Classic Bluetooth
  RFCOMM and BLE L2CAP, with detection reported at each separation.}
  \label{fig:transit}
\end{figure}

\begin{figure}[t]
  \centering
  \includegraphics[width=0.58\columnwidth]{fig7_club.png}
  \caption{The \emph{Club} tab --- the gamification surface where integrity is
  enforced. It reports mesh contribution, the scoring rule and the session
  history.}
  \label{fig:club}
\end{figure}

%---------------------------------------------------------------------------
% IV. METHOD
%---------------------------------------------------------------------------
\section{Cyber and Gamification Method}
\label{sec:method}

\subsection{Threat model}
We model a mesh of $N$ nodes in which $k$ are Byzantine. A Byzantine node can
(i) inject forged beacons --- spoofed SSIDs, rogue APs or phantom BLE
advertising --- to manufacture artificial disturbances, and (ii) publish forged
crossing reports (crossings it never observed) to inflate its own reward or
shift the fused estimate. We assume no trusted third party and no requirement
that nodes hold valid certificates, since the sensing radios are unmodified
commodity hardware.

\subsection{Byzantine influence on the mesh mean}
Let $C$ be the set of crossing reports and let the published count be the
thresholded consensus $C^* = f(\{c_j\})$ where each $c_j$ is a node's estimate
and $f$ down-weights reports far from the consensus. A Byzantine report
$c_k = c_{\mathrm{true}} + \Delta$ moves the consensus toward $\Delta$ in
proportion to its weight. The bounded-influence property of the mesh is the
key measurable outcome: over 160 audited empty scenes the mesh produced
\emph{zero} false confirmations, demonstrating that the agreement tolerance
bounds the influence of any single contributor.

\subsection{The alignment argument}
We formalise the reward as $R_i = R_{\mathrm{base}}(x_i) + \lambda \cdot
s_i \cdot g_i$, where $s_i$ is a reliability score inversely related to the
distance between node $i$'s estimate and the mesh mean, $g_i$ is a fusion
multiplier from multi-node confirmation, and $\lambda$ weights the bonus. A
Byzantine node that publishes $\Delta \neq 0$ maximises its base reward but
reduces $s_i$, so the marginal benefit of deception is negative for
$\lambda$ sufficiently large. In the limit, $R_i$ is maximised at
$s_i = 1$, which occurs only when the node reports the consensus value. The
incentive and the security objective therefore coincide.

\begin{algorithm}[t]
\caption{Reputation-weighted mesh confirmation}
\label{alg:mesh}
\begin{algorithmic}[1]
\Require set of node reports $\{c_j\}$, tolerance $\tau$, max distance $D$
\Ensure published count $C^*$, reliability scores $\{s_i\}$
\State Compute robust consensus $\mu = \mathrm{median}(\{c_j\})$
\For{each report $c_j$}
  \State $s_j \gets \max(0, 1 - |c_j - \mu| / D)$
\EndFor
\State Down-weight reports with $|c_j - \mu| > \tau$
\State $C^* \gets$ thresholded, down-weighted consensus
\State Award reliability bonus $\lambda\, s_j\, g_j$ per accepted report
\State \Return $C^*, \{s_j\}$
\end{algorithmic}
\end{algorithm}

\subsection{Future cyber defence: physical-layer anchoring}
Our planned defence replaces identifier-based trust with physical-layer
fingerprints~\cite{zhang2025survey,zhang2023wifi}. A beacon advertising a
spoofed SSID cannot replicate the intrinsic hardware impairments of a genuine
transmitter, so DFA-fused and multi-domain preamble features provide an
anti-spoofing layer~\cite{dfa2025}. Combined with channel-feature authentication
and environment-semantics-aware hardening~\cite{esaNet}, this converts the
reliability bonus into a security primitive rooted in physics rather than in
shared secrets. The two-node transit counter can similarly be validated using
round-trip-time analysis to detect rogue APs that spoof link identities.

%---------------------------------------------------------------------------
% V. EVALUATION
%---------------------------------------------------------------------------
\section{Evaluation}

\subsection{Harness}
Scenes are replayed through the production Kotlin classes, so a candidate
change must win on complete scenes, not row-level aggregates.
Table~\ref{tab:corpus} summarises the corpus.

\begin{table}[t]
  \centering
  \caption{Evaluation corpus replayed through the production engine.}
  \label{tab:corpus}
  \footnotesize
  \begin{tabular}{@{}lccc@{}}
    \toprule
    Experiment & Scenes & Purpose \\
    \midrule
    Multi-band fusion & 1,512 & per-crossing entity typing \\
    Occupied scenes per band & $378 \times 3$ & per-band classifier yield \\
    Ranging (near / far 55~m APs) & 288 & dBm--metre gate, band validation \\
    Epoch delta & 120 & movement vs.\ false positives \\
    Empty-scene audit & 160 & false positives, false confirmations \\
    \bottomrule
  \end{tabular}
\end{table}

\subsection{Fusion and per-band accuracy}
Across 1,512 scenes the classifier types a crossing correctly on BLE in
57.7\,\% of cases against 25.7\,\% (Wi-Fi 5~GHz) and 16.1\,\% (Wi-Fi 2.4~GHz);
equal-weight fusion was rejected for dragging BLE to 31\,\%
(Table~\ref{tab:bands}).

\begin{table}[t]
  \centering
  \caption{Per-band reference constants and measured per-crossing accuracy.}
  \label{tab:bands}
  \footnotesize
  \begin{tabular}{@{}lccccr@{}}
    \toprule
    Band & $n$ & $R_{1\,\mathrm{m}}$ & $A_0$ & Per-crossing & Max leg \\
    \midrule
    BLE (2.4~GHz) & 2.4 & $-42$~dBm & 7.2~dB & \textbf{57.7\,\%} & 24~m \\
    Wi-Fi 2.4~GHz & 2.6 & $-38$~dBm & 9.1~dB & 16.1\,\% & 58~m \\
    Wi-Fi 5~GHz & 2.9 & $-35$~dBm & 11.5~dB & 25.7\,\% & 68~m \\
    \bottomrule
  \end{tabular}
\end{table}

\subsection{Ranging and the gate}
On 288 scenes with near (10~m) and far (55~m) AP layouts, naive log-distance
inversion produced 14--23~m illusions on weak links. The $-75$~dBm gate and the
$f \in [0.15, 0.92]$ band suppress the illusion entirely, and links at the
$-90$~dBm floor are excluded from accuracy mathematics. The result is a range
that is sometimes refused, but never silently wrong --- an explicit honesty
property.

\subsection{Scan epochs}
Replacing naive last-sample comparison with the 20-sample median collapsed the
false-positive rate from 25--62\,\% to \textbf{zero at every tested distance}
(Table~\ref{tab:epoch}), demonstrating that the epoch rule is highly specific
even on empty scenes --- the property that makes it resistant to fabricated
disturbances.

\begin{table}[t]
  \centering
  \caption{Movement detection and false-positive rate for the epoch rule.}
  \label{tab:epoch}
  \footnotesize
  \begin{tabular}{@{}lcc@{}}
    \toprule
    Distance & Movement detected & False positives (naive) \\
    \midrule
    2~m & 100\,\% & 25\,\% \\
    4~m & 88\,\% & 34\,\% \\
    6~m & 81\,\% & 41\,\% \\
    8~m & 75\,\% & 48\,\% \\
    12~m & 62\,\% & 55\,\% \\
    20~m & 44\,\% & 62\,\% \\
    \bottomrule
  \end{tabular}
\end{table}

\subsection{Mesh and transit}
Over Classic Bluetooth RFCOMM and BLE L2CAP the hybrid classifier achieves
100\,\% fused detection and 100\,\% entity typing at 3--10~m and 94\,\% at
15~m; unconfirmed dips are held below threshold. With 2--5 cooperating
instances the mesh confirms events with zero false confirmations across 160
empty scenes and extends the useful radius from 4--12~m toward 20--40~m.

\subsection{Reliability bonus as trust}
The gamification layer estimates an explicit fused accuracy and rewards
participation with a per-crossing XP schedule, a $\times 2$ fusion bonus and
$+15$~XP per mesh-confirmed crossing, and a $\times 1.25$ reliability
multiplier for devices nearest the mesh mean. Because the reward derives from
the same fused estimate that is published, a device cannot inflate its score by
reporting noise. The calibration record shows the evidence rejected equal-weight
fusion and naive last-sample delta, both of which would have been exploitable
under the Byzantine model.

\subsection{Condition matrix}
The harness varies entity lateral profile, occlusion depth, closest approach,
AP layout, link occupancy and scene emptiness, bounding the regimes in which
each figure holds and the settings over which the integrity argument applies.

%---------------------------------------------------------------------------
% VI. CONCLUSIONS
%---------------------------------------------------------------------------
\section{Conclusions and Future Work}
MeshCrowd demonstrates that device-free sensing can be joined with
gamification such that the incentive and the security objective coincide. The
system fuses multi-band features into entity typing, converts measured dBm into
a gated range with explicit refusal semantics, detects change between two
instants of the field without inventing movement, and counts transit across a
two-node Bluetooth barrier --- all while a reliability bonus rewards the devices
nearest the fused estimate. Because the reward is computed from the estimate it
incentivises, reporting honesty is the rational strategy, and the mesh mean is
harder to attack than a system that rewards raw volume.

\textbf{Future work} proceeds on three fronts. First, a \emph{physical-layer
fingerprinting} module that anchors trust in hardware impairments rather than
identifiers~\cite{zhang2025survey,zhang2023wifi,dfa2025}, validated on the same
falsification protocol, with environment-semantics-aware
hardening~\cite{esaNet}. Second, \emph{reputation-weighted fusion} treating the
reliability bonus as an auditable, published score so that the "club" becomes a
transparent trust ledger. Third, \emph{longitudinal gamification evaluation}
with control groups, following the methodological critique of the gamified
security literature~\cite{gwenhure2025}, to measure whether the aligned reward
actually changes behaviour rather than only self-reported intention.

\textbf{Summary.} The central contribution is a demonstration that the two
problems --- making a sensing mesh trustworthy and keeping a voluntary mesh
alive --- can be solved by the same mechanism. When the reward is a function of
agreement with the consensus, the incentive to cheat is also the signal that
identifies the cheat.

%---------------------------------------------------------------------------
% REFERENCES
%---------------------------------------------------------------------------
\begin{thebibliography}{21}

\bibitem{wilson2010} J.~Wilson and N.~Patwari, ``Radio tomographic imaging with
wireless networks,'' \emph{IEEE Trans.\ Mobile Comput.}, vol.~9, no.~5,
pp.~621--632, May 2010, doi: 10.1109/TMC.2010.8.

\bibitem{januszkiewicz2018} {\L}.\,Januszkiewicz, ``Analysis of human body
shadowing effect on wireless sensor networks operating in the 2.4~GHz band,''
\emph{Sensors}, vol.~18, no.~10, art.~3412, Oct. 2018, doi: 10.3390/s18103412.

\bibitem{li2024} Z.~Li, A.~Dubey, S.~Shen, N.~K.~Kundu, J.~Rao, and
R.~Murch, ``Radio tomographic imaging with reconfigurable intelligent surfaces
for integrated sensing and communication,'' \emph{IEEE Trans.\ Wireless
Commun.}, vol.~23, no.~11, pp.~15784--15797, Nov. 2024, doi:
10.1109/TWC.2024.3433011.

\bibitem{ramirez2021} R.~Ramirez, C.-Y.~Huang, C.-A.~Liao, P.-T.~Lin, H.-W.~Lin,
and S.-H.~Liang, ``A practice of BLE RSSI measurement for indoor positioning,''
\emph{Sensors}, vol.~21, no.~15, art.~5181, Jul. 2021, doi: 10.3390/s21155181.

\bibitem{milano2024} F.~Milano, H.~da Rocha, M.~Laracca, L.~Ferrigno,
A.~Esp\'iritu Santo, J.~Salvado, and V.~Paciello, ``BLE-based indoor
localization: Analysis of some solutions for performance improvement,''
\emph{Sensors}, vol.~24, no.~2, art.~376, Jan. 2024, doi: 10.3390/s24020376.

\bibitem{bai2020} L.~Bai, F.~Ciravegna, R.~Bond, and M.~Mulvenna, ``A low cost
indoor positioning system using Bluetooth Low Energy,'' \emph{IEEE Access},
vol.~8, pp.~136858--136871, 2020, doi: 10.1109/ACCESS.2020.3012342.

\bibitem{khan2023} D.~Khan and I.~W.-H.~Ho, ``CrossCount: Efficient device-free
crowd counting by leveraging transfer learning,'' \emph{IEEE Internet Things
J.}, vol.~10, no.~5, pp.~4049--4058, Mar. 2023.

\bibitem{zhou2020} R.~Zhou, Z.~Gong, X.~Lu, and Y.~Fu, ``WiFlowCount:
Device-free people flow counting by exploiting Doppler effect in commodity
WiFi,'' \emph{IEEE Syst.~J.}, vol.~14, no.~4, pp.~4919--4930, Dec. 2020.

\bibitem{dmip2023} ``Non-intrusive people counting and identification
simultaneously with commodity WiFi devices,'' in \emph{Proc.\ 6th Int.\ Conf.\
Digital Medicine and Image Processing (DMIP '23)}, 2023, pp.~62--68, doi:
10.1145/3637684.3637694.

\bibitem{zhang2025survey} J.~Zhang, F.~Ardizzon, M.~Piana, G.~Shen, and
S.~Tomasin, ``Physical layer-based device fingerprinting for wireless security:
From theory to practice,'' \emph{IEEE Trans.\ Inf.\ Forensics Security},
2025 (to appear).

\bibitem{zhang2023wifi} J.~Zhang, Z.~Xu, J.~Li, Q.~Dai, Z.~Ling, and M.~Yang,
``Wi-Fi device identification based on multi-domain physical layer
fingerprint,'' \emph{Comput.\ Commun.}, 2023, doi: 10.1016/j.comcom.2023.03.024.

\bibitem{dfa2025} ``Radio frequency fingerprinting for WiFi authentication based
on detrended fluctuation analysis,'' \emph{IET Inf.\ Secur.}, 2025, doi:
10.1049/ise2/8683522.

\bibitem{phyfinatt} J.~Huang, B.~Liu, C.~Miao, X.~Zhang, J.~Liu, L.~Su, Z.~Liu,
and Y.~Gu, ``PhyFinAtt: An undetectable attack framework against PHY layer
fingerprint-based WiFi authentication,'' \emph{IEEE Trans.\ Mobile Comput.},
vol.~23, no.~7, pp.~7753--7770, 2023.

\bibitem{esaNet} N.~Gao, Q.~Huang, C.~Li, S.~Jin, and M.~Matthaiou, ``EsaNet:
Environment semantics enabled physical layer authentication,'' \emph{IEEE
Wireless Commun.\ Lett.}, vol.~13, no.~1, Jan. 2024.

\bibitem{zhang2022} J.~Zhang, X.~Li, Z.~Shi, \emph{et al.}, ``A reputation-based
and privacy-preserving incentive scheme for mobile crowd sensing: A deep
reinforcement learning approach,'' \emph{Wirel.\ Netw.}, 2023, doi:
10.1007/s11276-022-03111-9.

\bibitem{wang2018} J.~Wang, M.~Li, Y.~He, H.~Li, K.~Xiao, and C.~Wang, ``A
blockchain based privacy-preserving incentive mechanism in crowdsensing
applications,'' \emph{IEEE Access}, vol.~6, pp.~17545--17556, 2018, doi:
10.1109/ACCESS.2018.2815365.

\bibitem{gwenhure2025} T.~Gwenhure \emph{et al.}, ``Gamified cybersecurity
initiatives: The trend, limits and implications for practice and research,''
\emph{J. Inf.\ Technol.\ Educ.\ Res.}, vol.~24, 2025.

\bibitem{horst2024gamify} A.~Horst \emph{et al.}, ``Gamification in workforce
training: Improving employees' self-efficacy and information security and data
protection behaviours,'' \emph{J.~Bus.\ Res.}, 2024, doi:
10.1016/j.jbusres.2024.114747.

\bibitem{cm2026gamify} ``Applying gamification to promote privacy and
information security awareness in the public sector,'' \emph{Comput.\ Secur.},
2026.

\end{thebibliography}

\end{document}
